\documentclass[11pt,a4paper]{amsart}
\usepackage[a4paper,margin=20mm]{geometry}
\usepackage{amsmath,amssymb,amsthm,graphicx,xcolor}
\usepackage[T1]{fontenc}
\usepackage{lmodern}
\usepackage[no-math]{fontspec}
\setmainfont{lmroman10-regular.otf}[BoldFont=lmroman10-bold.otf,ItalicFont=lmroman10-italic.otf,BoldItalicFont=lmroman10-bolditalic.otf]

\usepackage[numbers]{natbib}
\renewcommand\bibsection{\section*{References}}
\usepackage[colorlinks=true,linkcolor=violet,citecolor=violet,urlcolor=teal]{hyperref}
\urlstyle{rm}
\usepackage{pgfplots,stmaryrd,multirow,tabularx,makecell,microtype,paralist,booktabs,adjustbox}
\makeatletter\renewcommand\@seccntformat[1]{\csname the#1\endcsname\quad}\makeatother
\usepackage[compact]{titlesec}
\usepackage[textsize=tiny,disable]{todonotes}
\usepackage{mathpartir,cleveref,versions}
\excludeversion{exclproof}
\excludeversion{technical}
\newtheorem{theorem}{Theorem}
\newtheorem{lemma}[theorem]{Lemma}
\newtheorem{proposition}[theorem]{Proposition}
\newtheorem{corollary}[theorem]{Corollary}
\theoremstyle{definition}
\newtheorem{definition}[theorem]{Definition}
\newtheorem{example}[theorem]{Example}
\newtheorem{remark}[theorem]{Remark}
\newtheorem{claim}[theorem]{Claim}
\AtBeginDocument{\def\shorttitle{Pushdown normal-form bisimulation}}
\renewenvironment{proof}[1][Proof]{\par\noindent\textit{#1.}\space}{\par}

\makeatletter
%%%%%% -- vk.sty
%%%%%% Usefull general macros
%%%%%% Last modified: Jan 24, 2022 - 15:33:16

%\RequirePackage[usenames,dvipsnames]{color} % this is for the \todo macro
\RequirePackage{alltt}
\RequirePackage{xspace}
\RequirePackage{xargs}

%%%%%%%%%%%%%%%%%%%%%%%%%%%%%%%
%%% Mathematics Environment %%%
%%%%%%%%%%%%%%%%%%%%%%%%%%%%%%%

% better ensuremath (from mmm.sty)
\newcommand{\ma}[1]{\ensuremath{#1}\xspace}

% (inductive) proof cases
\newcommand{\icase}[1]{\text{\textbf{Case~}{#1}\textbf{:}}}
%{{\flushleft\textbf{Case}~#1:}}

% modulo with nicer spacing
\def\imod#1{\allowbreak\mkern10mu({\operator@font mod}\,\,#1)}

% text operators
\def\tiff{\ma{\mathrel{\text{iff}}}}
\def\tand{\ma{\mathrel{\text{and}}}}
\def\timplies{\ma{\mathrel{\mathit{implies}}}}
\def\where{\ma{\mathrel{\vert}}}
\def\wehave{\ma{.~}}
\def\st{\ma{.~}}

% definition equal
\newcommand{\defeq}{\ma{\mathrel{\stackrel{%
        \mathrm{\text{\tiny def}}}{=}}}}

%%%%%%%%%%%%%%%
%%% Spacing %%%
%%%%%%%%%%%%%%%

% basic space length unit = 1 math unit = 0.0555em
\newlength{\basicspacelen}
\newlength{\basictablength}
\setlength{\basicspacelen}{0.0555em}
\setlength{\basictablength}{1em}

% variable width space
\newcommand{\vsp}[1][1]{\hspace*{#1\basicspacelen}}
% larger space
\newcommand{\vtab}[1][1]{\hspace*{#1\basictablength}}
% operator space (from mmm.sty)
%\newcommand{\opsp}[1][1]{\mathop{}{}}
\newcommand{\opsp}[1][1]{\vsp[#1]}

%%%%%%%%%%%%%
%%% Boxes %%%
%%%%%%%%%%%%%
% group boxz
\newcommand{\nbox}[2][t]{\ensuremath{\nboxstart{#1}#2\nboxstop}}
\newcommand{\nboxstart}[1]{\begin{array}[#1]{@{}l@{}}}
\newcommand{\nboxstop}{\end{array}}
\newcommand{\nsplit}[1][1]{\\ \vtab[#1]}
% align boxz
\newcommand{\abox}[2][t]{\ensuremath{\aboxstart{#1}#2\aboxstop}}
\newcommand{\aboxstart}[1]{\begin{array}[#1]{@{}r@{~}l@{}}}
\newcommand{\aboxstop}{\end{array}}
\newcommand{\asplit}[1][1]{\\ \vtab[#1]}
% new align boxz
\newcommandx{\Abox}[9][1=t,2=~,3=l,4=l,5=l,6=l,7=l,8=l]
  {\ma{\Aboxstart{#1}{#2}{#3}{#4}{#5}{#6}{#7}{#8}{l}#9\Aboxstop}}
\newcommand{\Aboxstart}[9]
  {\begin{array}[#1]{@{}#3@{#2}#4@{#2}#5@{#2}#6@{#2}#7@{#2}#8@{#2}#9@{}}}
\newcommand{\Aboxstop}{\end{array}}
\newcommand{\Asplit}[1][1]{\\\vtab[#1]}
% new align boxz
\newcommandx{\Cbox}[9][1=t,2=~,3=c,4=c,5=c,6=c,7=c,8=c]
  {\ma{\Aboxstart{#1}{#2}{#3}{#4}{#5}{#6}{#7}{#8}{l}#9\Aboxstop}}
% align both columns to the center
\newcommand{\ccbox}[1]{\ensuremath{\ccboxstart#1\ccboxstop}}
\newcommand{\ccboxstart}{\begin{array}[t]{@{}c@{\quad}c@{}}}
\newcommand{\ccboxstop}{\end{array}}
\newcommand{\ccsplit}[1]{\\}
% align both columns to the left
\newcommand{\lbox}[1]{\ensuremath{\lboxstart#1\lboxstop}}
\newcommand{\lboxstart}{\begin{array}[t]{@{}l@{\vsp}l@{}}}
\newcommand{\lboxstop}{\end{array}}
% centered box
\newcommand{\cbox}[2][t]{\ensuremath{\cboxstart{#1}#2\cboxstop}}
\newcommand{\cboxstart}[1]{\begin{array}[#1]{@{}c@{}}}
\newcommand{\cboxstop}{\end{array}}
\newcommand{\csplit}[1][1]{\\ \vtab[#1]}
% 3 collumn box
\newcommand{\threebox}[1]{\ensuremath{\threeboxstart#1\threeboxstop}}
\newcommand{\threeboxstart}{\begin{array}[t]{@{}r@{~}l@{~}l@{}}}
\newcommand{\threeboxstop}{\end{array}}
% new box, vertically alligned at the bottom
\newcommand{\bbox}[1]{\ensuremath{\bboxstart#1\bboxstop}}
\newcommand{\bboxstart}{\begin{array}[b]{@{}c@{}}}
\newcommand{\bboxstop}{\end{array}}
\newcommand{\bsplit}[1]{\\}



%%%%%%%%%%%%%%%%%%%%%%%%%
%%% Mathematical Sets %%%
%%%%%%%%%%%%%%%%%%%%%%%%%

%%% Symbols for the sets of the natural, real, rational,
%%% and complex numbers

\def\N{{I\mkern-4.7mu N}}
\def\R{{\textrm{I\mkern-3mu R}}}
\def\Q{{\@QC Q}}
\def\C{{\@QC C}}
\def\@QC#1{\mathpalette{\setbox0=\hbox\bgroup$\rm}%
  {\egroup C$\egroup\rm\rlap{\kern0.4\wd0\vrule
  width 0.05\wd0 height 0.97\ht0 depth -0.01\ht0}%
  #1\bgroup}}

\def\sdif{\ma{\mathbin{\backslash}}}
%%%%%%%%%%%%%%
%%% Floats %%%
%%%%%%%%%%%%%%
% Turn off the stupid acmconf rule above the caption
%\@setflag \@caprule = \@false

\def\dstart{\hbox to \hsize{\vrule depth 0pt\hrulefill\vrule depth 0pt}}
\def\dend{\hbox to \hsize{\phantom{\vrule height 5pt}\hrulefill\vrule height 0pt}}

\newcommand{\Hline}{%
\vskip\medskipamount % or other desired dimension
\leaders\vrule width \textwidth\vskip0.4pt % or other desired thickness
\vskip\medskipamount % ditto
\nointerlineskip
}

%\newlength{\myfigure@old@abovedisplayskip}
%\newlength{\myfigure@old@belowdisplayskip}
\newenvironment{myfigure*}[2]
{\begin{figure*}[t]
 \def\figtitle{#1}
 \def\figlabel{#2}
% \setlength{\myfigure@old@abovedisplayskip}{\abovedisplayskip}
% \setlength{\myfigure@old@belowdisplayskip}{\belowdisplayskip}
% \setlength{\abovedisplayskip}{0pt}
% \setlength{\belowdisplayskip}{0pt}
 \hrule
 %\fontsize{8.5}{9}\selectfont%
 %\footnotesize\selectfont
}{\hrule
 \caption{\figtitle\label{\figlabel}}
% \setlength{\abovedisplayskip}{\myfigure@old@abovedisplayskip}
% \setlength{\belowdisplayskip}{\myfigure@old@belowdisplayskip}
\end{figure*}}

\newenvironment{myfigure}[2]
{\begin{figure}[t]
 \def\figtitle{#1}%
 \def\figlabel{#2}%
% \hrule\vspace{-2ex}    %%%%MH
\hrule
}{
% \vspace{-2ex}\hrule    %%%%MH
\hrule
 \caption{\figtitle\label{\figlabel}}%
 \end{figure}%
}
\newenvironment{myfigureopt}[3]
{\begin{figure}[#1]
 \def\figtitle{#2}%
 \def\figlabel{#3}%
 \hrule
 \footnotesize\selectfont
}{
 \hrule
 \caption{\figtitle\label{\figlabel}}%
 \end{figure}%
}



% Left-justified, boxed titles in floats:
\newlength{\@fltitlelen}
\newlength{\@fltitlelentmp}
\newcommand{\floattitle}[1]
  {\settowidth{\@fltitlelentmp}{~}%
   \setlength{\@fltitlelen}{\textwidth}%
   \addtolength{\@fltitlelen}{-\@fltitlelentmp}%
   ~\makebox[\@fltitlelen][l]{\fbox{~{#1}~}}}
\newcommand{\floattitleb}[1]
  {~\fbox{~{#1}~}}


%%% Conditional to test for empty arguments.
%%% It doesn't work in some weird cases when \iftrue or \iffalse is used, but
%%% I find this macro more standard; see source.
%%% source: http://www.ctan.org/pub/tex-archive/info/aro-bend/answer.002

%    \@car is actually already defined in latex.tex, but for
%    maximum robustness it needs to have the \long prefix:
\long\def\@car#1#2\@nil{#1}
\long\def\@first#1#2{#1}
\long\def\@second#1#2{#2}
\long\def\ifempty#1{\expandafter\ifx\@car#1@\@nil @\@empty
  \expandafter\@first\else\expandafter\@second\fi}



%%%%%%%%%%%%%%%%%%%%%%%
%%% Inference rules %%%
%%%%%%%%%%%%%%%%%%%%%%%

% inference rule
% usage: \irule[name][irefname]{premises}{conclusion}
% usage: \irule*[name][irefname]{premises}{conclusion}

\newcommand{\irule}{\@ifnextchar*\irule@top\irule@second}
\newcommand{\irule@second}{\@ifnextchar!\irule@singleline\irule@side}

\def\irule@top*{\irule@top@prime}
\newcommandx{\irule@top@prime}[4][1=,2=]{
  \ma{\inferrule[{#1}]{#3}{#4}}%
  \ifempty{#2}{}%
   {\expandafter\gdef\csname ir@#2\endcsname{\ma{\RefTirName{#1}}}}%
}
\def\irule@side{\irule@side@prime}
\newcommandx{\irule@side@prime}[4][1=,2=]{
  \mprset{center}\ma{\inferrule{#3}{#4}\;\RefTirName{#1}}%
  \ifempty{#2}{}%
   {\expandafter\gdef\csname ir@#2\endcsname{\ma{\RefTirName{#1}}}}%
}
\def\irule@singleline!{\irule@singleline@prime}
\newcommandx{\irule@singleline@prime}[4][1=,2=]{
  \ma{\ifempty{#1}{}{\RefTirName{#1}: }  #4 & \ifempty{#3}{&}{\text{if~} & #3}}%
  \ifempty{#2}{}%
   {\expandafter\gdef\csname ir@#2\endcsname{\ma{\RefTirName{#1}}}}%
}

% Reference name of inference rules
\newcommand{\iref}[1]{\ma{%
%  \def\tmp@name{\@irefname{#1}}%
  \vkifundef{ir@#1}%
    \@questionmark%
  \else%
    \csname ir@#1\endcsname%
  \fi%
}}
\newcommand{\@questionmark}{\textbf{?}\message{Warning: there were undefined references to rules!}}
%\newcommand{\@irefname}[1]{\csname ir@#1\endcsname}

%%%%%%%%%%%%%%%%%%%%
%%% Conditionals %%%
%%%%%%%%%%%%%%%%%%%%

% if undefined macro. Usage: \vkifundef{A} ... \else ... \fi
\def\vkifundef#1{\expandafter\ifx\csname#1\endcsname\relax}

%%% Conditional to test for empty arguments.
%%% It doesn't work in some weird cases when \iftrue or \iffalse is used, but
%%% I find this macro more standard; see source.
%%% source: http://www.ctan.org/pub/tex-archive/info/aro-bend/answer.002

%    \@car is actually already defined in latex.tex, but for
%    maximum robustness it needs to have the \long prefix:
\long\def\@car#1#2\@nil{#1}
\long\def\@first#1#2{#1}
\long\def\@second#1#2{#2}
\long\def\ifempty#1{\expandafter\ifx\@car#1@\@nil @\@empty
  \expandafter\@first\else\expandafter\@second\fi}


%%%%%%%%%%%%%%%%%%%%%%%%%%%%%%%%%
%%% 'Work in progress' macros %%%
%%%%%%%%%%%%%%%%%%%%%%%%%%%%%%%%%

% draft heading
\def\Draft{\RequirePackage{fancyhdr}\chead{Draft, %
\today}}

% to do macros
%%\def\pretodo#1{\typeout{*** TODO: #1}}
%%\def\todo#1{\pretodo{#1}\text{{{\color{red}\bfseries TODO:} #1}}}

%\newcommand{\todo}[1][]{{\color{red}\bfseries TODO:} #1}

%%% Symbol for the set of integer numbers

%You can try to make a symbol for the integer numbers from
%two Z's as above, but that won't look very nice. Instead
%you can make your own symbol using Metafont

%\DeclareFontFamily{U}{cmmz}{}
%\DeclareFontShape{U}{cmmz}{m}{n}{<-8><8><9><10>%
%  <10.95><12><14.4><17.28><20.74><24.88>cmmz10}{}
%\DeclareSymbolFont{mz}{U}{cmmz}{m}{n}
%\DeclareMathSymbol\Z{\mathord}{mz}{"5A}

% bold math font!
%\def\bmath#1{\mbox{\boldmath$#1$}}
\def\bmath#1{\mathchoice{\mbox{\boldmath$\displaystyle#1$}}
{\mbox{\boldmath$\textstyle#1$}}
{\mbox{\boldmath$\scriptstyle#1$}}
{\mbox{\boldmath$\scriptscriptstyle#1$}}}

\newcommand{\freshfrom}{\ma{\mathrel{\sharp}}}
\makeatother

\makeatletter

\RequirePackage{mathtools}
\RequirePackage{amssymb}
\RequirePackage{bbold}

%\RequirePackage{courier}
\RequirePackage[scaled]{beramono}
%\RequirePackage{inconsolata}
%\RequirePackage{anyfontsize}
\RequirePackage{listings}
%\usepackage{DejaVuSansMono}
%\lstset{language=ML,basicstyle=\footnotesize\sffamily,mathescape}
%\lstset{language=ML,basicstyle=\fontsize{9}{10}\ttfamily,mathescape}
%\lstset{language=ML,basicstyle=\fontsize{9}{10}\rmfamily,mathescape}
\lstset{language=ML,basicstyle=\linespread{0.8}\footnotesize\ttfamily,mathescape}
% the language name
\newcommand{\lang}{\ma{\lambda^{\textsf{imp}}}}
\lstset{
 morekeywords={ref}
}
% tool names
\newcommand\tool[1]{{\sc#1}}
\newcommand\Hobbit{\tool{Hobbit}}
\newcommand\SLHobbit{\tool{PDNF-Bisim}}


% some language syntax
\newcommand{\true}{\ma{\mathsf{tt}}}
\newcommand{\false}{\ma{\mathsf{ff}}}
\newcommand{\nil}{\ma{\mathsf{()}}}
\newcommand{\lam}[2][f]{\ensuremath{\mathsf{fix} #1 (#2) .}}
\newcommand{\app}[1]{\ma{#1 \,}}
\newcommand{\asgn}{\ma{\mathbin{:=}}}
\newcommand{\ifk}{\ma{\mathop{\mathsf{if}}}}
\newcommand{\thenk}{\ma{\mathop{\mathsf{then}}}}
\newcommand{\elsek}{\ma{\mathop{\mathsf{else}}}}
\newcommand{\cond}[3]{\ma{\ifk #1 \thenk #2 \elsek #3}}
\newcommand{\letk}{\ma{\mathop{\mathsf{let}}}}
\newcommand{\ink}{\ma{\mathrel{\mathsf{in}}}}
\newcommand{\elet}[3]{\ma{\letk #1 \mathbin{=} #2 \ink #3}}
\newcommand{\new}[2]{\ma{\mathsf{ref}\, #1 \mathbin{=} #2 \mathop{\mathsf{in}}}}
\newcommand{\deref}[1]{\ma{{!}#1}}
\newcommand{\arithop}[2][]{\ma{\mathop{\ifempty{#1}{op}{#1}}(#2)}}

\newcommand{\typing}[4]{\ma{#1; #2 \mathrel{\vdash} #3 : #4}}
% contexts
\newcommand{\hole}[1][\cdot]{\ma{[#1]}}
% Semantics
\newcommand{\redconf}[2]{\ma{\langle #1 \mathop{;} #2 \rangle}} % reduction configuration
\newcommand{\red}{\ma{\rightarrow}}
\newcommand{\redbase}{\ma{\hookrightarrow}}
\newcommand{\funred}{\ma{\succ}}
% LTS
\newcommand{\conf}[5]{\ma{\langle #1 \mathop{;} #2 \mathop{;} #3 \mathop{;} #4 \mathop{;} #5\rangle}}
\newcommand{\oldpconf}[5]
{\ma{\color{blue}\langle\color{black} #1 \mathop{;} #2 \mathop{;} #3 \mathop{;} #4 \mathop{;} #5 \color{blue}\rangle\color{black}}
}
\newcommand{\oldoconf}[5]
{\ma{{\color{red}\langle} #1 \mathop{;} #2 \mathop{;} {#3} \mathop{;} #4  \mathop{;} #5 {\color{red}\rangle}}
}
\newcommand{\pconf}[4]
{\ma{\color{blue}\llparenthesis\color{black} #1 \mathop{;} #2 \mathop{;} #3 \mathop{;} #4 \color{blue}\rrparenthesis\color{black}}
}
\newcommand{\oconf}[4]
{\ma{\color{red}\llparenthesis\color{black} #1 \mathop{;} #2 \mathop{;} {#3} \mathop{;} #4 \color{red}\rrparenthesis\color{black}}
}
\newcommand{\botconf}{\ma{\langle \bot \rangle}}
\newcommand{\trans}[2][]{\ma{\ifempty{#1}{\xrightarrow{#2}}{\xrightarrow[#1]{#2}}}}
\newcommand{\transinv}[2][]{\ma{\ifempty{#1}{\xleftarrow{#2}}{\xrightarrow[#1]{#2}}}}
\newcommand{\transs}[2][]{\ma{\ifempty{#1}{\xrightarrow{#2}}{\xrightarrow[#1]{#2}}\hspace{-.7ex}{}^*}}
\newcommand{\transsinv}[2][]{\ma{\ifempty{#1}{\xleftarrow{#2}}{\xrightarrow[#1]{#2}}\hspace{-.7ex}{}^*}}
\newcommand{\wtrans}[2][]{\ma{\ifempty{#1}{\xRightarrow{#2}}{\xRightarrow[#1]{#2}}}}
\newcommand{\betatrans}{\realtrans{\tau}_{\bmath{\beta}}}
\newcommand{\betatransinv}{\transinv{\tau}_{\bmath{\beta}}}
\newcommand{\betatranss}{\ma{\mathrel{\realtrans{\tau}\hspace{-.7ex}{}_{\bmath{\beta}}^*}}}
\newcommand{\betatranssinv}{\ma{\mathrel{\xleftarrow{\tau}\hspace{-.7ex}{}_{\bmath{\beta}}^*}}}
\newcommand\satrans[1]{\trans[{\sf SAT}]{#1}}
\newcommand\satranss[1]{\transs[{\sf SAT}]{#1}}
\newcommand{\noe}{\ma{\cdot}}
\newcommand{\emptyA}{\ma{\cdot}}
\newcommand{\emptyK}{\ma{\cdot}}
\newcommand{\emptyG}{\ma{\cdot}}
\newcommand{\emptyS}{\ma{\cdot}}
\newcommand{\emptylist}{\ma{\epsilon}}

% ultimate pattern matching
\newcommand{\ulpatt}{\ma{\mathop{\mathsf{ulpatt}}}}


% LTS labels
\newmuskip\uappskip
\DeclareMathOperator{\OpApp}{%
  \mkern\uappskip
  \underline{\mkern-\uappskip \mathsf{\color{red}call\color{black}}\mkern-\uappskip}
  \mkern\uappskip
}
\DeclareMathOperator{\OpRet}{%
  \mkern\uappskip
  \underline{\mkern-\uappskip \mathsf{\color{red}ret\color{black}}\mkern-\uappskip}
  \mkern\uappskip
}
\DeclareMathOperator{\PropApp}{%
  \mathsf{\color{blue}call\color{black}}}
%  \mkern\uappskip
%  \overline{\mkern-\uappskip \textsf{app}\mkern-\uappskip}
%  \mkern\uappskip
%}
\DeclareMathOperator{\PropRet}{%
  \mathsf{\color{blue}ret\color{black}}}
%  \mkern\uappskip
%  \overline{\mkern-\uappskip \textsf{ret}\mkern-\uappskip}
%  \mkern\uappskip
%}
\newcommand{\lopapp}[2]{\ma{\OpApp(#1, #2)}}
\newcommand{\lopret}[1]{\ma{\OpRet(#1)}}
\newcommandx{\lpropapp}[2]{\ma{\PropApp(#1,#2)}}
\newcommand{\lpropret}[1]{\ma{\PropRet(#1)}}
\newcommand{\lterm}{\ma{\downarrow}}

% some basic sets
\newcommand{\Typ}{\ma{\mathsf{Type}}} % Types
\newcommand{\Bool}{\ma{\mathsf{bool}}} % Bool type
\newcommand{\Int}{\ma{\mathsf{int}}} % Int type
\newcommand{\Unit}{\ma{\mathsf{unit}}} % Unit type
\newcommand{\arrow}{\ma{\rightarrow}} % Arrow type
\newcommand{\Var}{\ma{\mathsf{Var}}} % Variables
\newcommand{\Con}{\ma{\mathsf{Const}}} % Constants
\newcommand{\Exp}{\ma{\mathsf{Exp}}} % Expressions
\newcommand{\Val}{\ma{\mathsf{Val}}} % Values
\newcommand{\EC}{\ma{\mathsf{ECxt}}} % Evaluation Contexts
\newcommand{\Cxt}{\ma{\mathsf{Cxt}}} %  Contexts
\newcommand{\Loc}{\ma{\mathsf{Loc}}} % set of locations
\newcommand{\St}{\ma{\mathsf{St}}} % set of stores
\newcommand{\Conf}{\ma{\mathsf{Conf}}} % set of configurations (for operational semantics)
\newcommand{\Abs}{\ma{\mathbb{A}}} % set of abstract names
\newcommand{\Nat}{\ma{\mathbb{N}}} % set of abstract names

% macros for typesetting math
\newcommand{\mis}{\ma{\mathbin{::=}}}
\newcommand{\mor}{\ma{\mathrel{\vert}}}
\newcommand{\morcondensed}{\ma{\hspace{1.5pt}{\vert}\hspace{1.5pt}}}

% meta-operations
\newcommand{\sub}[2]{\ma{[#2/#1]}} % substitution
\newcommand{\subnr}[2]{\ma{[#2/#1]_{\mathsf{nr}}}} % name-reuse substitution
\newcommand{\stupd}[2]{\ma{[#1\mapsto#2]}} % store extension/update
\newcommand{\concr}[2]{\ma{\{#2/#1\}}} % store extension/update
\newcommand{\trm}{\ma{\Downarrow}}
\newcommand{\dvg}{\ma{\Uparrow}}

\newcommand{\dom}[1]{\ma{\mathsf{dom}(#1)}} % domain
\newcommand{\rng}[1]{\ma{\mathsf{rng}(#1)}} % range
\newcommand{\nxt}[1]{\ma{\mathsf{nxt}(#1)}} % next abstract value

\newcommand{\fv}[1]{\ma{\mathsf{fv}(#1)}} % free variables
\newcommand{\fl}[1]{\ma{\mathsf{fl}(#1)}} % free locations
\newcommand{\an}[1]{\ma{\mathsf{an}(#1)}} % abstract names
\newcommand{\supp}[1]{\ma{\mathsf{supp}(#1)}} % support of a permutation
\newcommand{\fran}[1]{\ma{\mathsf{fran}(#1)}} % fresh abstract names on a transition
%stackless:
\newcommand{\gext}[2][]{\ma{\mathsf{gext}_{#1}(#2)}} % fresh abstract names on a transition

% permutations
\newcommand{\perm}[2]{\ma{(#1\leftrightarrow #2)}}
\newcommand{\pic}[1][\pi]{#1\cdot}

% Environments
\newcommand{\maps}[2]{\ma{\prescript{#1}{}{#2}}}
\newcommand{\sizeof}[1]{\ma{\mathopen{|}#1\mathclose{|}}}
\newcommandx{\sepconj}[2][1=,2=]{\ma{\mathrel{{\oplus}_{#1}^{#2}}}}
\renewcommandx{\interleave}[1][1=]{\ma{\mathbin{\#_{#1}}}}

% RELATIONS
\newcommand\fresh{\ma{\mathrel{\#}}}
\newcommand{\cxteq}{\ma{\mathrel{\equiv}}}
\newcommand{\betaeq}{\ma{\mathrel{\simeq_\beta}}}

\newcommand{\bisim}{\@ifstar{\bisimRel}{\bisimNoRel}}
\newcommand{\bisimRel}[1]{\ma{\mathrel{\bisimNoRel{#1}}}}
\newcommand{\bisimNoRel}[1]{\ma{\mathcal{#1}}}

\newcommand{\bisimil}{\ma{\approx}}
\newcommand{\sbisimil}{\ma{\sim}}
\newcommand{\simil}{\ma{\mathrel{\let\scriptstyle\textstyle\substack{\sqsubset\\[-0.1ex]\approx}}}}
\newcommand{\ssimil}{\ma{\mathrel{\let\scriptstyle\textstyle\substack{\sqsubset\\[-0.2ex]\sim}}}}

% up-to techniques
% Macros that become mathrel when adding a star.
% first argument is the name of the function
% second argument is optional, allowing for an exponent
% third argument is the argument

\newcommandx{\Relation}[4][2=,3=]{\ma{\mathrel{\NotRelation{#1}[#2][#3]{#4}}}}
\newcommandx{\NotRelation}[4][2=,3=]{\ma{\mathsf{#1}\ifempty{#2}{}{^{#2}}\ifempty{#3}{}{_{#3}}\ifempty{#4}{}{(#4)}}}
\newcommand{\RelationOrNot}[1]{\@ifstar{\Relation{#1}}{\NotRelation{#1}}}

\newcommand{\ut}{\@ifstar{\Relation}{\NotRelation}} % general up to with an additional argument after the star

\newcommand{\utbeta}{\RelationOrNot{beta}} % up to beta reductions
\newcommand{\utperm}{\RelationOrNot{perm}}% up to permutations
\newcommand{\utcxt}{\RelationOrNot{cxt}}% up to context
\newcommand{\utfoldv}{\RelationOrNot{fldv}} % up to fold
\newcommand{\utfolde}{\RelationOrNot{flde}} % up to fold
%\newcommand{\utan}{\RelationOrNot{an}} % up to abstract name sets
\newcommand{\weakrel}{\prec^{\mathsf{wk}}}
\newcommand{\weakrelinv}{\succ^{\mathsf{wk}}}
\newcommand{\utweak}[1][]{\RelationOrNot{weak\ma{{}_{#1}}}}% up to weakening Gammas
\newcommand{\utsepconj}[1][]{\RelationOrNot{sep_{#1}}}% up to separating conjunction
\newcommand{\utgext}{\RelationOrNot{gext}}% up to separating conjunction
\newcommand{\utId}{\RelationOrNot{id}} % up to identity
\newcommand{\utRefl}{\RelationOrNot{\constf{refl}}} % up to identity
\newcommand{\utunreachable}{\RelationOrNot{{unr}}} % up to unreachable alpha
\newcommand{\asympgc}{\asymp^{\mathsf{gc}}}
\newcommand{\gcrel}{\prec^{\mathsf{gc}}}
\newcommand{\gcrelinv}{\succ^{\mathsf{gc}}}
\newcommand{\utgc}{\RelationOrNot{{gc}}} % up to unreachable alpha
\newcommand{\utreentry}{\RelationOrNot{{reent}}} % up to proponent function reentry
\newcommand{\utabs}{\RelationOrNot{{abs}}} % up to abstraction
\newcommand{\uttaut}{\RelationOrNot{{taut}}} % up to tautology
\newcommand{\constf}[1]{\ma{{\mathsf{c}_{#1}}}}
\newcommand{\pow}[1]{\ma{\mathcal{P}(#1)}}
\newcommand{\utnr}{\RelationOrNot{nr}}% up to weakening Gammas
\newcommand{\utnrinv}{\RelationOrNot{nr\ma{^{-1}}}}% up to weakening Gammas

% machinery for proving up to techniques due to Damien Pous
\newcommand{\WP}{\ma{\mathbf{wp}}}
\newcommand{\progress}[1][]{\ma{\mathrel{\overset{#1}{\rightsquigarrow}}}}
\newcommand{\gfp}[1]{\ma{\mathsf{gfp}(#1)}}
\newcommand{\companion}[1][]{\ma{\mathbf{t}_{#1}}}

% extras
%\newcommand{\defqed}{%
%  \renewcommand{\qedsymbol}{$\lozenge$}%
%  \pushQED{\qed}%
%  \qedhere%
%  \popQED%
%}
\newcommand{\defqed}{}
\renewcommand{\qedhere}{}

\newcommand{\bij}{\ma{\mathbin{\leftrightarrow}}}
\newcommand{\comp}{\,{{\circ}}\,}
\let\uplusorig=\uplus
%\def\uplus{\hspace{1pt}{\uplusorig}\hspace{1pt}}
\def\uplus{\mskip2mu{\uplusorig}\mskip2mu}
  
%%%%%%%%%%%%%%%%
%%% Comments %%%
%%%%%%%%%%%%%%%%

\newcommand\commentbox[2]{{\color{#1}#2}}
\newcommand\nt[1]{\commentbox{cyan}{#1}}
\newcommand\vk[1]{\commentbox{blue}{#1}}
\newcommand\yy[1]{\commentbox{red}{#1}}

%%%%%%%%%%%%%%%%%%%%%%
%%% Symbolic Stuff %%%
%%%%%%%%%%%%%%%%%%%%%%
\newcommand{\symval}[1]{\ma{\widetilde{#1}}}%symbolic values

% SETS 
\newcommand{\SymCon}{\ma{\mathsf{SymCon}}} % Sym Constants
\newcommand{\SymEnv}{\ma{\mathsf{SymEnv}}} % set of stores
\newcommand{\Prop}{\ma{\mathsf{Prop}}} % set of stores
\newcommand{\Assign}{\ma{\mathsf{Assign}}} % set of stores

% SYMBOLIC COMPONENTS
\newcommand{\sig}{\ma{\sigma}} % sym env
\newcommand{\modelM}{\ma{\delta}} % model

% SYMBOLIC CONFIGS
\newcommand{\symredconf}[3]{\ma{ #1 \vdash \langle #2 \mathop{;} #3 \rangle}} % reduction configuration
\newcommand{\symltsconf}[2]{\ma{ #1 \vdash #2}} % sig |- C
%\newcommand{\symconf}[6]{\ma{ #1 \vdash \conf{#2}{#3}{#4}{#5}{#6}}} % symbolic lts conf

% RELATIONS
\newcommand{\symbisim}[2]{\ma{\bisim{#1}_{#2}}} % sym bisim
\newcommand{\symeq}[3]{\ma{#1 \stackrel{#2}{=} #3}} % equivalence in model

% META
%\newcommand{\phiof}[1]{\ma{{#1}^\circ}}%symbolic values
\newcommand{\concretise}[2]{\ma{#1\{#2\}}} % replace in #1 all things mapped in #2
\newcommand{\symbisimil}[1]{\ma{\stackrel{#1}{\bisimil}}}



%%%%%%%%%%%%%%
% TIKZ STUFF %
%%%%%%%%%%%%%%
\usepackage{tikz}
\usetikzlibrary{positioning}
\usetikzlibrary{automata,shapes.multipart}

\tikzstyle{onode} = [circle, minimum width=1cm,text centered, draw=black, fill=red!30, align=center]
\tikzstyle{pnode} = [circle, minimum width=1cm,text centered, draw=black, fill=blue!30, align=center]

\tikzstyle{ostate} = [
		align=center,
        state,
        circle split, 
        minimum width=1cm,text centered, draw=black, fill=red!30]
\tikzstyle{pstate} = [
		align=center,
        state,
        circle split, 
        minimum width=1cm,text centered, draw=black, fill=blue!30]

\usetikzlibrary{
    decorations.markings,
    arrows.meta  %% <--
}

\tikzset{  
	strike through/.style={
    postaction=decorate,
    decoration={
      markings,
      mark=at position 0.5 with {
	        \draw[-] (-6pt,-6pt) -- (6pt, 6pt);
      		}
    	}
	}
}

\newcommand{\papp}[2]{\lpropapp{#1}{#2}}
\newcommand{\pret}[1]{\textsf{ret}\langle#1\rangle}
\newcommand{\oapp}[2]{\underline{\textsf{app}}\langle#1,#2\rangle}
\newcommand{\oret}[1]{\underline{\textsf{ret}}\langle#1\rangle}


%%% Pair LTS %%%

\newcommand{\bbisim}{\@ifstar{\bbisimRel}{\bbisimNoRel}}
\newcommand{\bbisimRel}[1]{\ma{\mathrel{\bbisimNoRel{#1}}}}
\newcommand{\bbisimNoRel}[1]{\ma{\mathbb{#1}}}
\newcommand{\bconf}[2]{\ma{\bigl\langle\!\!\bigl\langle #1 \mathop{;} #2 \bigr\rangle\!\!\bigr\rangle}}
\newcommand{\thickbar}[1]{\mathbf{\widehat{\text{\ma{#1}}}}}
\newcommand{\co}[1]{\ma{\thickbar{#1}}}

%% bold symbols for pair LTS
\DeclareSymbolFont{bbsymbol}{U}{bbold}{m}{n}
\DeclareMathSymbol{\bbeta}{\mathord}{bbsymbol}{"0C}
\DeclareMathSymbol{\bSigma}{\mathalpha}{bbsymbol}{"06}
\DeclareMathSymbol{\bC}{\mathalpha}{bbsymbol}{`C}

\newcommand{\bbisimil}{\ma{\hat\approx}}
\newcommand{\bsimil}{\ma{\mathrel{\let\scriptstyle\textstyle\substack{\widehat\sqsubset\\[-0.1ex]\approx}}}}

\newcommand{\bmerge}[1]{\ma{\llbracket #1 \rrbracket}}
\makeatother
\crefname{section}{Sec.}{Sec.(s)}
\crefname{appendix}{Appx.}{Appx.(s)}
\crefname{figure}{Fig.}{Fig.(s)}
\crefname{theorem}{Thm.}{Thm.(s)}
\crefname{definition}{Def.}{Def.(s)}
\crefname{proposition}{Prop.}{Prop.(s)}
\crefname{corollary}{Cor.}{Cor.(s)}
\crefname{lemma}{Lem.}{Lem.(s)}
\crefname{example}{Ex.}{Ex.(s)}
\crefname{remark}{Rem.}{Rem.(s)}
\crefname{equation}{}{}

\newcommand\cutout[1]{}
\AtBeginDocument{\special{pdf:mapline fvmr8r  BeraSansMono-Roman     " TeXBase1Encoding ReEncodeFont " <8r.enc <fvmr8a.pfb}}
\AtBeginDocument{\special{pdf:mapline fvmro8r BeraSansMono-Oblique   " TeXBase1Encoding ReEncodeFont " <8r.enc <fvmro8a.pfb}}
\AtBeginDocument{\special{pdf:mapline fvmb8r  BeraSansMono-Bold      " TeXBase1Encoding ReEncodeFont " <8r.enc <fvmb8a.pfb}}
\AtBeginDocument{\special{pdf:mapline fvmbo8r BeraSansMono-BoldOb    " TeXBase1Encoding ReEncodeFont " <8r.enc <fvmbo8a.pfb}}
\AtBeginDocument{\special{pdf:mapline bbold17 BBOLD10 <bbold10.pfb}}
\AtBeginDocument{\special{pdf:mapline bbold12 BBOLD10 <bbold10.pfb}}
\AtBeginDocument{\special{pdf:mapline bbold10 BBOLD10 <bbold10.pfb}}
\AtBeginDocument{\special{pdf:mapline bbold9 BBOLD10 <bbold10.pfb}}
\AtBeginDocument{\special{pdf:mapline bbold8 BBOLD7 <bbold7.pfb}}
\AtBeginDocument{\special{pdf:mapline bbold7 BBOLD7 <bbold7.pfb}}
\AtBeginDocument{\special{pdf:mapline bbold6 BBOLD5 <bbold5.pfb}}
\AtBeginDocument{\special{pdf:mapline bbold5 BBOLD5 <bbold5.pfb}}

\renewcommand\ttfamily{\fontencoding{T1}\fontfamily{fvm}\selectfont}
